% =============================================================================
% APPENDIX PC: FORMAL-PCT: Parameter Context Transformation
% =============================================================================
% LANGUAGE: english
% =============================================================================
%
% HEADER BLOCK
% -----------------------------------------------------------------------------
% Appendix:      PC
% Category:      FORMAL (Formalization)
% Version:       1.0 (February 2026)
% Dependencies:  V (CORE-WHEN), BBB (CORE-WHERE), AN (METHOD-LLMMC)
% Primary Ch.:   9 (Context as Endogenous)
% Secondary Ch.: 4 (Empirical Foundations), 5 (Complementarity), 8 (Mathematical)
% -----------------------------------------------------------------------------
%
% This appendix formalizes the central theoretical contribution of EBF:
% Behavioral parameters are not constants but context-dependent functions.
% θ = f(Ψ, 10C) - The Parameter Context Transformation (PCT) framework.
%
% =============================================================================

% =============================================================================
% CROSS-REFERENCE STRUCTURE
% =============================================================================
%
%                    ┌─────────────┐
%                    │  Chapter 9  │ (WHEN - Context)
%                    │   CORE-V    │
%                    └──────┬──────┘
%                           │
%          ┌────────────────┼────────────────┐
%          ↓                ↓                ↓
%   ┌─────────────┐  ┌─────────────┐  ┌─────────────┐
%   │  Appendix V │  │ Appendix PC │  │Appendix BBB │
%   │  CORE-WHEN  │◄─│ FORMAL-PCT  │─►│ CORE-WHERE  │
%   │  (Context)  │  │   (THIS)    │  │ (Parameters)│
%   └─────────────┘  └──────┬──────┘  └─────────────┘
%                           │
%          ┌────────────────┼────────────────┐
%          ↓                ↓                ↓
%   ┌─────────────┐  ┌─────────────┐  ┌─────────────┐
%   │ Appendix AN │  │  Paper-YAML │  │Appendix CAL │
%   │ METHOD-LLMMC│  │    (SSOT)   │  │ METHOD-CAL  │
%   │  (Priors)   │  │             │  │(Calibration)│
%   └─────────────┘  └─────────────┘  └─────────────┘
%
% =============================================================================

\chapter{PC FORMAL-PCT: Parameter Context Transformation}
\label{app:formal-pct}

% =============================================================================
% ABSTRACT
% =============================================================================

\begin{abstract}
This appendix formalizes the central theoretical contribution of EBF:
behavioral parameters are context-dependent functions, not universal constants.
We introduce the Parameter Context Transformation (PCT) framework with the
central equation $\theta_B = \theta_A \times \prod_i M(\Delta\Psi_i) \times \prod_j N(\Delta 10C_j)$,
enabling systematic parameter transfer across contexts. The framework treats
cross-study variation as signal (information about context-dependence) rather
than noise (estimation error). We provide axioms, worked examples, and an
iterative learning paradigm where each new paper improves all future predictions.
\end{abstract}

% =============================================================================
% QUICK REFERENCE BOX
% =============================================================================

\begin{tcolorbox}[colback=blue!5!white, colframe=blue!75!black, title=Quick Reference: Parameter Context Transformation]
\textbf{Central Equation (PCT):}
\[
\theta(B) = \theta(A) \times T(\Delta\Psi, \Delta 10C)
\]

\textbf{Multiplicative Form:}
\[
\theta_B = \theta_A \times \prod_i M(\Delta\Psi_i) \times \prod_j N(\Delta 10C_j)
\]

\textbf{Context Function (Goal):}
\[
\theta = F_\theta(\Psi, 10C) = \theta_{base} \times \prod_i (1 + \beta_i \cdot \Psi_i) \times \prod_j (1 + \gamma_j \cdot 10C_j)
\]

\textbf{Key Symbols:}
\begin{itemize}
    \item $\theta$ = Behavioral parameter (e.g., $\lambda$, $\beta$, $\alpha$, $\mu$)
    \item $\Psi$ = Context vector $(\Psi_I, \Psi_S, \Psi_C, \Psi_K, \Psi_F, \Psi_T, \Psi_M, \Psi_E)$
    \item $10C$ = Utility dimension vector (WHO, WHAT, HOW, WHEN, WHERE, ...)
    \item $M(\Delta\Psi_i)$ = Context multiplier for dimension $i$
    \item $N(\Delta 10C_j)$ = Utility dimension multiplier for dimension $j$
\end{itemize}
\end{tcolorbox}

% =============================================================================
% APPENDIX SCOPE BOX
% =============================================================================

\begin{tcolorbox}[colback=gray!5!white, colframe=gray!75!black, title=Appendix Scope]
\textbf{Ziel:} Formalize the theoretical claim that behavioral parameters are
context-dependent functions, and provide a methodology for parameter transformation.

\textbf{In-Scope:}
\begin{itemize}
    \item PCT axioms and formal definitions
    \item Transformation equation derivation
    \item Measurement contexts architecture
    \item Iterative learning paradigm
    \item Scientific novelty documentation
\end{itemize}

\textbf{Out-of-Scope:}
\begin{itemize}
    \item Specific parameter estimates (see Appendix BBB)
    \item Context dimension definitions (see Appendix V)
    \item LLMMC methodology (see Appendix AN)
\end{itemize}

\textbf{Constraints:}
\begin{itemize}
    \item All parameter values must trace to Paper-YAML (SSOT)
    \item Transformations must be falsifiable
\end{itemize}
\end{tcolorbox}

% =============================================================================
% FUNDAMENTAL QUESTION
% =============================================================================

\begin{tcolorbox}[colback=fadarkblue!10!white, colframe=fadarkblue,
    title=\textbf{Fundamental Question}]
\large
\textbf{How can behavioral parameters measured in context $A$ be systematically
transformed to predict behavior in context $B$?}
\normalsize

\vspace{0.5em}
Traditional behavioral economics treats parameters as universal constants
(e.g., $\lambda = 2.25$). The EBF recognizes that the same individual exhibits
different ``loss aversion'' in different contexts. PCT provides the formal
apparatus to model, estimate, and predict these context-dependent transformations.
\end{tcolorbox}


% =============================================================================
% SECTION 1: THE PARADIGM SHIFT
% =============================================================================

\section{The Paradigm Shift: From Constants to Functions}
\label{sec:pct-paradigm}

\subsection{The Traditional View}

In traditional behavioral economics, parameters are treated as universal constants:

\begin{equation}
\lambda = 2.25 \quad \text{(Loss Aversion)}
\end{equation}

\begin{equation}
\beta = 0.70 \quad \text{(Present Bias)}
\end{equation}

\begin{equation}
\alpha = 0.85 \quad \text{(Inequity Aversion)}
\end{equation}

This view implies that loss aversion ``is'' 2.25---a fixed property of human psychology.

\subsection{The Replication Problem}

Empirically, parameters show substantial variation across studies:

\begin{center}
\begin{tabular}{lcc}
\toprule
\textbf{Parameter} & \textbf{Traditional Value} & \textbf{Observed Range} \\
\midrule
Loss Aversion $\lambda$ & 2.25 & $[1.5, 5.0]$ \\
Present Bias $\beta$ & 0.70 & $[0.50, 0.95]$ \\
Inequity Aversion $\alpha$ & 0.85 & $[0.40, 1.20]$ \\
\bottomrule
\end{tabular}
\end{center}

\textbf{Traditional interpretation:} This variation is ``noise''---estimation error,
methodological differences, or sampling variation to be controlled for.

\subsection{The EBF Interpretation}

\begin{tcolorbox}[colback=green!5!white, colframe=green!75!black, title=Core Insight]
\textbf{The variation is not noise---it is the signal!}

Parameter variation across studies reflects systematic context-dependence:
\[
\theta = f(\Psi, 10C)
\]

Parameters are functions of context ($\Psi$) and utility dimensions ($10C$),
not universal constants.
\end{tcolorbox}

\begin{axiom}[PCT-1: Contextual Parameters]
\label{ax:pct-1}
Every behavioral parameter $\theta$ is a function of context $\Psi$ and
utility dimensions $10C$:
\[
\theta: \Psi \times 10C \to \mathbb{R}^+
\]
The ``constant'' reported in any study is $\theta(\Psi_{\text{study}}, 10C_{\text{study}})$---the
parameter value evaluated at that study's specific context.
\end{axiom}


% =============================================================================
% SECTION 2: THE TRANSFORMATION EQUATION
% =============================================================================

\section{The Parameter Context Transformation Equation}
\label{sec:pct-equation}

\subsection{The Central Equation}

Given a parameter $\theta$ measured in context $A$, we can predict its value
in context $B$ using the transformation equation:

\begin{equation}
\boxed{
\theta(\Psi_B, 10C_B) = \theta(\Psi_A, 10C_A) \times T(\Delta\Psi, \Delta 10C)
}
\label{eq:pct-central}
\end{equation}

where:
\begin{itemize}
    \item $\theta(\Psi_A, 10C_A)$ = Parameter measured in context $A$ (anchor)
    \item $\Delta\Psi = \Psi_B - \Psi_A$ = Context difference vector
    \item $\Delta 10C = 10C_B - 10C_A$ = Utility dimension difference vector
    \item $T(\cdot)$ = Transformation function
\end{itemize}

\subsection{Multiplicative Form}

For practical application, we decompose $T$ multiplicatively:

\begin{equation}
\theta_B = \theta_A \times \prod_{i=1}^{8} M(\Delta\Psi_i) \times \prod_{j=1}^{9} N(\Delta 10C_j)
\label{eq:pct-multiplicative}
\end{equation}

where:
\begin{itemize}
    \item $M(\Delta\Psi_i)$ = Multiplier for context dimension $i$ (Institutional, Social, Cognitive, ...)
    \item $N(\Delta 10C_j)$ = Multiplier for utility dimension $j$ (WHO, WHAT, HOW, ...)
\end{itemize}

\begin{axiom}[PCT-2: Multiplicative Decomposition]
\label{ax:pct-2}
Context effects are multiplicative and separable across dimensions:
\[
T(\Delta\Psi, \Delta 10C) = \prod_i M_i(\Delta\Psi_i) \times \prod_j N_j(\Delta 10C_j)
\]
This implies that context dimensions affect parameters independently (no interactions
in first-order approximation).
\end{axiom}

\subsection{The Context Multiplier Table}

Empirical estimates of $M(\Delta\Psi)$ from cross-study analysis:

\begin{center}
\begin{tabular}{llcc}
\toprule
\textbf{$\Psi$-Dimension} & \textbf{Direction} & \textbf{Multiplier} & \textbf{Example} \\
\midrule
$\Psi_S$ (Social) & Stigma $\downarrow$ & 0.80--0.90 & welfare $\to$ workplace \\
$\Psi_S$ (Social) & Stigma $\uparrow$ & 1.10--1.25 & private $\to$ public \\
$\Psi_I$ (Institutional) & Formal $\downarrow$ & 0.90--0.98 & bureaucratic $\to$ informal \\
$\Psi_I$ (Institutional) & Formal $\uparrow$ & 1.02--1.15 & informal $\to$ regulated \\
$\Psi_C$ (Cognitive) & Load $\uparrow$ & 1.05--1.20 & relaxed $\to$ stressed \\
$\Psi_C$ (Cognitive) & Load $\downarrow$ & 0.85--0.95 & stressed $\to$ relaxed \\
$\Psi_K$ (Cultural) & Norm shift & 0.70--1.40 & culture-dependent \\
$\Psi_F$ (Physical) & Digital $\to$ F2F & 1.10--1.30 & more personal \\
\bottomrule
\end{tabular}
\end{center}


% =============================================================================
% SECTION 3: THE CONTEXT FUNCTION
% =============================================================================

\section{Parameter Context Functions}
\label{sec:pct-functions}

\subsection{From Transformation to Function}

While PCT enables point-to-point transformation, the ultimate goal is to estimate
the full context function $F_\theta$:

\begin{equation}
\theta = F_\theta(\Psi, 10C)
\label{eq:context-function}
\end{equation}

\begin{axiom}[PCT-3: Estimable Functions]
\label{ax:pct-3}
For each behavioral parameter $\theta$, there exists an estimable function
$F_\theta: \mathbb{R}^{8} \times \mathbb{R}^{9} \to \mathbb{R}^+$ that maps
context and utility dimensions to parameter values.
\end{axiom}

\subsection{Functional Form}

We propose a multiplicative functional form:

\begin{equation}
F_\theta(\Psi, 10C) = \theta_{base} \times \prod_{i=1}^{8} (1 + \beta_i \cdot \Psi_i) \times \prod_{j=1}^{9} (1 + \gamma_j \cdot 10C_j)
\label{eq:context-function-form}
\end{equation}

where:
\begin{itemize}
    \item $\theta_{base}$ = Baseline parameter value (neutral context)
    \item $\beta_i$ = Sensitivity to context dimension $i$
    \item $\gamma_j$ = Sensitivity to utility dimension $j$
\end{itemize}

\subsection{Example: Loss Aversion Function}

\begin{tcolorbox}[colback=yellow!5!white, colframe=yellow!75!black, title=Hypothetical Estimate: $F_\lambda$]
\begin{equation}
\lambda(\Psi, 10C) = 2.00 \times (1 + 0.25 \cdot \text{Stigma}) \times (1 + 0.10 \cdot \text{Formal}) \times (1 + 0.15 \cdot \text{CogLoad}) \times \ldots
\end{equation}

\begin{center}
\begin{tabular}{lccc}
\toprule
\textbf{Coefficient} & \textbf{Estimate} & \textbf{SE} & \textbf{Interpretation} \\
\midrule
$\lambda_{base}$ & 2.00 & 0.15 & Baseline loss aversion \\
$\beta_S$ (Stigma) & 0.25 & 0.04 & +25\% per stigma level \\
$\beta_I$ (Formal) & 0.10 & 0.03 & +10\% per formality level \\
$\beta_C$ (CogLoad) & 0.15 & 0.05 & +15\% under cognitive load \\
$\beta_K$ (IndivNorm) & 0.20 & 0.06 & +20\% with individualism norm \\
$\gamma_{WHO}$ (Self) & 0.05 & 0.02 & +5\% when self (vs. other) \\
$\gamma_{AWARE}$ (Salience) & 0.12 & 0.04 & +12\% with high salience \\
\bottomrule
\end{tabular}
\end{center}

\textbf{Example predictions:}
\begin{align*}
\lambda(\text{welfare, self, salient}) &= 2.0 \times 1.25 \times 1.10 \times 1.15 \times 1.05 \times 1.12 = 3.74 \\
\lambda(\text{workplace, self, routine}) &= 2.0 \times 1.00 \times 1.05 \times 1.00 \times 1.05 \times 1.00 = 2.20
\end{align*}
\end{tcolorbox}


% =============================================================================
% SECTION 4: MEASUREMENT CONTEXTS ARCHITECTURE
% =============================================================================

\section{Measurement Contexts: The Data Architecture}
\label{sec:measurement-contexts}

\subsection{The SSOT Principle}

\begin{axiom}[PCT-4: Paper-YAML as SSOT]
\label{ax:pct-4}
Paper-YAML files are the Single Source of Truth (SSOT) for all parameter measurements.
Each \texttt{measurement\_context} documents:
\begin{enumerate}
    \item The parameter value estimate
    \item The $\Psi$-conditions under which it was measured
    \item The source location in the original paper
    \item The study type (theoretical, lab, field)
\end{enumerate}
\end{axiom}

\subsection{Schema Structure}

\begin{verbatim}
parameter_contributions:
  - symbol: "λ_R"
    ebf_id: "PAR-BEH-016"
    measurement_contexts:
      - context: "welfare_takeup"
        value_estimate: "high (λ_R ≈ 2.5)"
        psi_conditions:
          Ψ_I: "bureaucratic_application"
          Ψ_S: "welfare_stigma"
          Ψ_K: "self_reliance_norm"
        source_in_paper: "Section 6.2"
        study_type: "field_data"
        countries: ["USA"]
\end{verbatim}

\subsection{From Contexts to Triplets}

For estimation, we extract $(θ, Ψ, 10C)$ triplets from all Paper-YAMLs:

\begin{center}
\begin{tabular}{ccccccc}
\toprule
$\theta$ & $\Psi_I$ & $\Psi_S$ & $\Psi_K$ & WHO & $\ldots$ \\
\midrule
$\lambda = 2.5$ & bureaucratic & welfare\_stigma & self\_reliance & self & $\ldots$ \\
$\lambda = 1.8$ & professional & competence & team & self & $\ldots$ \\
$\lambda = 2.2$ & medical\_gate & illness\_stigma & privacy & self & $\ldots$ \\
$\vdots$ & $\vdots$ & $\vdots$ & $\vdots$ & $\vdots$ & $\ldots$ \\
\bottomrule
\end{tabular}
\end{center}


% =============================================================================
% SECTION 5: ITERATIVE LEARNING
% =============================================================================

\section{Iterative Learning: The Bayesian Foundation}
\label{sec:iterative-learning}

\subsection{The Learning Paradigm}

\begin{axiom}[PCT-5: Iterative Improvement]
\label{ax:pct-5}
Each new paper with \texttt{measurement\_contexts} improves the estimate of $F_\theta$
for all parameters. The improvement is monotonic: prediction error decreases with
more data points.
\end{axiom}

\subsection{Bayesian Update}

Formally, we update $F_\theta$ using Bayes' rule:

\begin{equation}
P(\beta | \text{Data}) \propto P(\text{Data} | \beta) \times P(\beta)
\label{eq:bayes-update}
\end{equation}

where:
\begin{itemize}
    \item $P(\beta)$ = Prior from LLMMC + existing literature
    \item $P(\text{Data} | \beta)$ = Likelihood of observed measurement contexts given coefficients
    \item $P(\beta | \text{Data})$ = Posterior after observing new paper
\end{itemize}

\subsection{Convergence}

\begin{tcolorbox}[colback=blue!5!white, colframe=blue!75!black, title=Convergence Guarantee]
Under mild regularity conditions:
\begin{enumerate}
    \item Posterior contracts to true parameter values
    \item Credible intervals become narrower
    \item Out-of-sample prediction error decreases monotonically
\end{enumerate}

With sufficient data, $F_\theta$ approximates the true context function.
\end{tcolorbox}

\subsection{Iteration Examples}

\begin{center}
\begin{tabular}{llccc}
\toprule
\textbf{Iteration} & \textbf{Data} & \textbf{Model} & \textbf{$R^2$} \\
\midrule
0 & N=0 (prior) & $F_\lambda = \lambda_{base} = 2.25$ & -- \\
1 & N$\approx$30 contexts & $F_\lambda(\Psi_S)$ & 0.25 \\
2 & N$\approx$150 contexts & $F_\lambda(\Psi_S, \Psi_I)$ & 0.45 \\
3 & N$\approx$500 contexts & $F_\lambda(\Psi, 10C)$ & 0.65--0.75 \\
\bottomrule
\end{tabular}
\end{center}


% =============================================================================
% SECTION 6: SCIENTIFIC NOVELTY
% =============================================================================

\section{Scientific Novelty}
\label{sec:novelty}

\subsection{What Exists in the Literature}

\begin{itemize}
    \item \textbf{Meta-analyses:} Report average parameter values across studies
    \item \textbf{Heterogeneity analyses:} Note ``large variance'' between studies
    \item \textbf{Cultural comparisons:} Hofstede dimensions, GLOBE, Inglehart
    \item \textbf{WEIRD critique:} Henrich et al.~show non-generalizability
    \item \textbf{Situationism debate:} Mischel, Ross \& Nisbett
\end{itemize}

\subsection{What Does NOT Exist (EBF Contribution)}

\begin{tcolorbox}[colback=red!5!white, colframe=red!75!black, title=Novel Contributions]
\begin{enumerate}
    \item \textbf{Formal specification:} $\theta = f(\Psi, 10C)$ with estimable coefficients
    \item \textbf{Systematic triplet collection:} $(\theta, \Psi, 10C)$ from papers via \texttt{measurement\_contexts}
    \item \textbf{Estimable context functions:} $F_\theta$ with coefficients $\beta_i$, $\gamma_j$
    \item \textbf{Paradigm:} ``Variation is signal, not noise''
    \item \textbf{Utility dimensions as moderators:} 10C integration
    \item \textbf{Transformation equation:} $\theta_B = \theta_A \times \prod_i M(\Delta\Psi_i)$
\end{enumerate}
\end{tcolorbox}

\subsection{Why This Did Not Exist Before}

\begin{enumerate}
    \item \textbf{Methodological focus:} ``Control for context'' instead of ``model through context''
    \item \textbf{Replication crisis framing:} Variation as problem to solve, not data to use
    \item \textbf{Missing data structure:} Papers don't report context systematically
    \item \textbf{Disciplinary silos:} Psychology (individual heterogeneity) vs. cultural economics (macro differences)
\end{enumerate}


% =============================================================================
% SECTION 7: WORKED EXAMPLE
% =============================================================================

\section{Worked Example: Rejection Sensitivity Transformation}
\label{sec:worked-example}

\subsection{Setup}

From Bénabou, Jaroszewicz \& Loewenstein (2022), we have two measurement contexts
for rejection sensitivity $\lambda_R$:

\begin{center}
\begin{tabular}{lll}
\toprule
\textbf{Context} & \textbf{$\lambda_R$} & \textbf{$\Psi$-Conditions} \\
\midrule
Welfare takeup & $\approx 2.5$ & $\Psi_I$: bureaucratic, $\Psi_S$: welfare\_stigma, $\Psi_K$: self\_reliance \\
Workplace help & $\approx 1.8$ & $\Psi_I$: professional, $\Psi_S$: competence\_signaling, $\Psi_C$: evaluation \\
\bottomrule
\end{tabular}
\end{center}

\subsection{Transformation Calculation}

\textbf{Step 1: Identify $\Delta\Psi$}

\begin{align*}
\Delta\Psi_S &: \text{welfare\_stigma} \to \text{competence\_signaling} = -0.3 \\
\Delta\Psi_I &: \text{bureaucratic} \to \text{professional} = -0.1 \\
\Delta\Psi_C &: \text{(none)} \to \text{performance\_evaluation} = +0.2
\end{align*}

\textbf{Step 2: Apply Multipliers}

\begin{align*}
M(\Delta\Psi_S = -0.3) &= 0.85 \quad \text{(less stigma $\to$ less sensitivity)} \\
M(\Delta\Psi_I = -0.1) &= 0.95 \quad \text{(less formal $\to$ slightly less)} \\
M(\Delta\Psi_C = +0.2) &= 1.10 \quad \text{(evaluation $\to$ more sensitivity)}
\end{align*}

\textbf{Step 3: Transform}

\begin{align*}
\lambda_R(\text{workplace}) &= \lambda_R(\text{welfare}) \times M(\Delta\Psi_S) \times M(\Delta\Psi_I) \times M(\Delta\Psi_C) \\
&= 2.5 \times 0.85 \times 0.95 \times 1.10 \\
&= 2.22
\end{align*}

\subsection{Comparison}

\begin{center}
\begin{tabular}{lcc}
\toprule
& \textbf{Measured} & \textbf{PCT Predicted} \\
\midrule
$\lambda_R$ (workplace) & 1.8 & 2.22 \\
\bottomrule
\end{tabular}
\end{center}

\textbf{Deviation:} 23\%---suggests additional factors not captured (e.g., $\Psi_F$: relationship type).


% =============================================================================
% SECTION 8: RESULTS
% =============================================================================

\section{Results}
\label{sec:pct-results}

\begin{theorem}[Parameter Context Dependence]
\label{thm:pct-main}
Under axioms PCT-1 through PCT-5, behavioral parameter $\theta$ measured in
context $A$ can be transformed to context $B$ via:
\[
\theta_B = \theta_A \times \prod_i M(\Delta\Psi_i) \times \prod_j N(\Delta 10C_j)
\]
with prediction error bounded by the estimation uncertainty of multipliers $M$ and $N$.
\end{theorem}

\begin{corollary}[Iterative Convergence]
\label{cor:convergence}
Under Bayesian updating with informative priors, the estimated context function
$\hat{F}_\theta$ converges to the true function $F_\theta$ as the number of
measurement contexts $n \to \infty$, with variance decreasing as $O(1/n)$.
\end{corollary}

\begin{result}[Cross-Study Variation]
The observed variation in replicated behavioral economics studies (e.g.,
$\lambda \in [1.5, 5.0]$) is not primarily estimation error but systematic
context-dependence that PCT can model and predict.
\end{result}


% =============================================================================
% SECTION 9: SUMMARY
% =============================================================================

\section{Summary}

\begin{tcolorbox}[colback=gray!5!white, colframe=gray!75!black, title=Summary Box]
\textbf{Core Claim:} Behavioral parameters are context-dependent functions, not constants.

\textbf{Central Equation:}
\[
\theta_B = \theta_A \times \prod_i M(\Delta\Psi_i) \times \prod_j N(\Delta 10C_j)
\]

\textbf{Research Program:} Estimate $F_\theta(\Psi, 10C)$ for each parameter using
measurement contexts from Paper-YAMLs.

\textbf{Iterative Learning:} Each new paper improves all future predictions.

\textbf{Scientific Contribution:} First formal framework for context-dependent
behavioral parameters with estimable coefficients.
\end{tcolorbox}


% =============================================================================
% GLOSSARY
% =============================================================================

\section{Glossary of Symbols}

\begin{center}
\begin{tabular}{ll}
\toprule
\textbf{Symbol} & \textbf{Definition} \\
\midrule
$\theta$ & Behavioral parameter (e.g., $\lambda$, $\beta$, $\alpha$) \\
$\Psi$ & Context vector with 8 dimensions \\
$\Psi_I$ & Institutional context dimension \\
$\Psi_S$ & Social context dimension \\
$\Psi_C$ & Cognitive context dimension \\
$\Psi_K$ & Cultural context dimension \\
$\Psi_F$ & Physical context dimension \\
$10C$ & Utility dimension vector (WHO, WHAT, HOW, WHEN, ...) \\
$F_\theta$ & Context function for parameter $\theta$ \\
$M(\Delta\Psi_i)$ & Context multiplier for dimension $i$ \\
$N(\Delta 10C_j)$ & Utility dimension multiplier for dimension $j$ \\
$\theta_{base}$ & Baseline parameter value (neutral context) \\
\bottomrule
\end{tabular}
\end{center}

For comprehensive symbol definitions, see Appendix GLS (Central Glossary).


% =============================================================================
% CRITICAL FOUNDATIONS
% =============================================================================

\section{Critical Foundations}
\label{sec:pct-foundations}

\begin{enumerate}
    \item \textbf{Prospect Theory (Kahneman \& Tversky, 1979):} Establishes that
    $\lambda > 1$ (loss aversion) is a fundamental behavioral parameter, but treats
    it as context-independent.

    \item \textbf{WEIRD Critique (Henrich et al., 2010):} Demonstrates that behavioral
    parameters vary systematically across populations, motivating context-dependence.

    \item \textbf{Situationist Psychology (Mischel, 1968; Ross \& Nisbett, 1991):}
    Shows that behavior is more situation-dependent than trait-dependent, providing
    psychological foundation for PCT.

    \item \textbf{Meta-Analytic Heterogeneity:} Systematic reviews showing large
    $I^2$ heterogeneity in behavioral parameters across studies---PCT reinterprets
    this as informative data.

    \item \textbf{Bayesian Learning (Gelman et al., 2013):} Provides the statistical
    machinery for iterative updating of context functions.
\end{enumerate}


% =============================================================================
% OPEN ISSUES
% =============================================================================

\section{Open Issues and Research Agenda}

\begin{enumerate}
    \item \textbf{Interaction terms:} When do context dimensions interact non-additively?
    \item \textbf{Nonlinear effects:} Are multipliers constant or context-dependent themselves?
    \item \textbf{Cultural variation:} How do $\beta_K$ coefficients vary across cultures?
    \item \textbf{Temporal dynamics:} How do context functions evolve over time?
    \item \textbf{Individual heterogeneity:} How to integrate individual-level variation?
\end{enumerate}


% =============================================================================
% REFERENCES
% =============================================================================

\section{References}

\nocite{bcm_master}

Key foundational papers:
\begin{itemize}
    \item Kahneman \& Tversky (1979): Prospect Theory
    \item Fehr \& Schmidt (1999): Inequity Aversion
    \item Laibson (1997): Quasi-Hyperbolic Discounting
    \item Henrich et al. (2010): WEIRD critique
    \item Bénabou, Jaroszewicz \& Loewenstein (2022): Fear of Asking (worked example)
\end{itemize}

See Master Bibliography (bcm\_master.bib) for full citations.


% =============================================================================
% CROSS-REFERENCE FOOTER
% =============================================================================

\vspace{1cm}
\hrule
\vspace{0.3cm}
\textbf{Cross-References:}
\begin{itemize}
    \item Appendix V (CORE-WHEN): Context dimensions $\Psi$
    \item Appendix BBB (CORE-WHERE): Parameter values and sources
    \item Appendix AN (METHOD-LLMMC): Prior generation methodology
    \item Appendix CAL (METHOD-CAL): Calibration and validation
    \item Appendix GLS: Central Glossary
\end{itemize}

\end{document}
